\documentclass[a4paper,12pt,twoside]{memoir}

% Castellano
\usepackage[spanish,es-tabla]{babel}
\selectlanguage{spanish}
\usepackage[utf8]{inputenc}
\usepackage[T1]{fontenc}
\usepackage{lmodern} % Scalable font
\usepackage{microtype}
\usepackage{placeins}
\usepackage{eurosym}

\RequirePackage{booktabs}
\RequirePackage[table]{xcolor}
\RequirePackage{xtab}
\RequirePackage{multirow}

% Links
\usepackage[colorlinks]{hyperref}
\hypersetup{
	allcolors = {red}
}

% Ecuaciones
\usepackage{amsmath}

% Rutas de fichero / paquete
\newcommand{\ruta}[1]{{\sffamily #1}}

% Párrafos
\nonzeroparskip
% Huérfanas y viudas
\widowpenalty100000
\clubpenalty100000

% Imagenes
\usepackage{graphicx}
\newcommand{\imagen}[2]{
	\begin{figure}
		\centering
		\includegraphics[width=0.9\textwidth]{#1}
		\caption{#2}\label{fig:#1}
	\end{figure}
	\FloatBarrier
}

\newcommand{\imagenPeque}[2]{
	\begin{figure}%[!h]
		\centering
		\includegraphics[width=0.4\textwidth]{#1}
		\caption{#2}\label{fig:#1}
	\end{figure}
	\FloatBarrier
}
\newcommand{\imagenPequeDos}[2]{
	\begin{figure}%[!h]
		\centering
		\includegraphics[width=0.6\textwidth]{#1}
		\caption{#2}\label{fig:#1}
	\end{figure}
	\FloatBarrier
}

\newcommand{\imagenflotante}[2]{
	\begin{figure}%[!h]
		\centering
		\includegraphics[width=0.9\textwidth]{#1}
		\caption{#2}\label{fig:#1}
	\end{figure}
}


% El comando \figura nos permite insertar figuras comodamente, y utilizando
% siempre el mismo formato. Los parametros son:
% 1 -> Porcentaje del ancho de página que ocupará la figura (de 0 a 1)
% 2 --> Fichero de la imagen
% 3 --> Texto a pie de imagen
% 4 --> Etiqueta (label) para referencias
% 5 --> Opciones que queramos pasarle al \includegraphics
% 6 --> Opciones de posicionamiento a pasarle a \begin{figure}
\newcommand{\figuraConPosicion}[6]{%
  \setlength{\anchoFloat}{#1\textwidth}%
  \addtolength{\anchoFloat}{-4\fboxsep}%
  \setlength{\anchoFigura}{\anchoFloat}%
  \begin{figure}[#6]
    \begin{center}%
      \Ovalbox{%
        \begin{minipage}{\anchoFloat}%
          \begin{center}%
            \includegraphics[width=\anchoFigura,#5]{#2}%
            \caption{#3}%
            \label{#4}%
          \end{center}%
        \end{minipage}
      }%
    \end{center}%
  \end{figure}%
}

%
% Comando para incluir imágenes en formato apaisado (sin marco).
\newcommand{\figuraApaisadaSinMarco}[5]{%
  \begin{figure}%
    \begin{center}%
    \includegraphics[angle=90,height=#1\textheight,#5]{#2}%
    \caption{#3}%
    \label{#4}%
    \end{center}%
  \end{figure}%
}
% Para las tablas
\newcommand{\otoprule}{\midrule [\heavyrulewidth]}
%
% Nuevo comando para tablas pequeñas (menos de una página).
\newcommand{\tablaSmall}[5]{%
 \begin{table}
  \begin{center}
   \rowcolors {2}{gray!35}{}
   \begin{tabular}{#2}
    \toprule
    #4
    \otoprule
    #5
    \bottomrule
   \end{tabular}
   \caption{#1}
   \label{tabla:#3}
  \end{center}
 \end{table}
}

%
% Nuevo comando para tablas pequeñas (menos de una página).
\newcommand{\tablaSmallSinColores}[5]{%
 \begin{table}[H]
  \begin{center}
   \begin{tabular}{#2}
    \toprule
    #4
    \otoprule
    #5
    \bottomrule
   \end{tabular}
   \caption{#1}
   \label{tabla:#3}
  \end{center}
 \end{table}
}

\newcommand{\tablaApaisadaSmall}[5]{%
\begin{landscape}
  \begin{table}
   \begin{center}
    \rowcolors {2}{gray!35}{}
    \begin{tabular}{#2}
     \toprule
     #4
     \otoprule
     #5
     \bottomrule
    \end{tabular}
    \caption{#1}
    \label{tabla:#3}
   \end{center}
  \end{table}
\end{landscape}
}

%
% Nuevo comando para tablas grandes con cabecera y filas alternas coloreadas en gris.
\newcommand{\tabla}[6]{%
  \begin{center}
    \tablefirsthead{
      \toprule
      #5
      \otoprule
    }
    \tablehead{
      \multicolumn{#3}{l}{\small\sl continúa desde la página anterior}\\
      \toprule
      #5
      \otoprule
    }
    \tabletail{
      \hline
      \multicolumn{#3}{r}{\small\sl continúa en la página siguiente}\\
    }
    \tablelasttail{
      \hline
    }
    \bottomcaption{#1}
    \rowcolors {2}{gray!35}{}
    \begin{xtabular}{#2}
      #6
      \bottomrule
    \end{xtabular}
    \label{tabla:#4}
  \end{center}
}

%
% Nuevo comando para tablas grandes con cabecera.
\newcommand{\tablaSinColores}[6]{%
  \begin{center}
    \tablefirsthead{
      \toprule
      #5
      \otoprule
    }
    \tablehead{
      \multicolumn{#3}{l}{\small\sl continúa desde la página anterior}\\
      \toprule
      #5
      \otoprule
    }
    \tabletail{
      \hline
      \multicolumn{#3}{r}{\small\sl continúa en la página siguiente}\\
    }
    \tablelasttail{
      \hline
    }
    \bottomcaption{#1}
    \begin{xtabular}{#2}
      #6
      \bottomrule
    \end{xtabular}
    \label{tabla:#4}
  \end{center}
}

%
% Nuevo comando para tablas grandes sin cabecera.
\newcommand{\tablaSinCabecera}[5]{%
  \begin{center}
    \tablefirsthead{
      \toprule
    }
    \tablehead{
      \multicolumn{#3}{l}{\small\sl continúa desde la página anterior}\\
      \hline
    }
    \tabletail{
      \hline
      \multicolumn{#3}{r}{\small\sl continúa en la página siguiente}\\
    }
    \tablelasttail{
      \hline
    }
    \bottomcaption{#1}
  \begin{xtabular}{#2}
    #5
   \bottomrule
  \end{xtabular}
  \label{tabla:#4}
  \end{center}
}



\definecolor{cgoLight}{HTML}{EEEEEE}
\definecolor{cgoExtralight}{HTML}{FFFFFF}

%
% Nuevo comando para tablas grandes sin cabecera.
\newcommand{\tablaSinCabeceraConBandas}[5]{%
  \begin{center}
    \tablefirsthead{
      \toprule
    }
    \tablehead{
      \multicolumn{#3}{l}{\small\sl continúa desde la página anterior}\\
      \hline
    }
    \tabletail{
      \hline
      \multicolumn{#3}{r}{\small\sl continúa en la página siguiente}\\
    }
    \tablelasttail{
      \hline
    }
    \bottomcaption{#1}
    \rowcolors[]{1}{cgoExtralight}{cgoLight}

  \begin{xtabular}{#2}
    #5
   \bottomrule
  \end{xtabular}
  \label{tabla:#4}
  \end{center}
}


















\graphicspath{ {./img/} }

% Capítulos
\chapterstyle{bianchi}
\newcommand{\capitulo}[2]{
	\setcounter{chapter}{#1}
	\setcounter{section}{0}
	\chapter*{#2}
	\addcontentsline{toc}{chapter}{#2}
	\markboth{#2}{#2}
}

% Apéndices
\renewcommand{\appendixname}{Apéndice}
\renewcommand*\cftappendixname{\appendixname}

\newcommand{\apendice}[1]{
	%\renewcommand{\thechapter}{A}
	\chapter{#1}
}

\renewcommand*\cftappendixname{\appendixname\ }

% Formato de portada
\makeatletter
\usepackage{xcolor}
\newcommand{\tutor}[1]{\def\@tutor{#1}}
\newcommand{\course}[1]{\def\@course{#1}}
\newcommand{\cotutor}[1]{\def\@cotutor{#1}}
\definecolor{cpardoBox}{HTML}{E6E6FF}
\def\maketitle{
  \null
  \thispagestyle{empty}
  % Cabecera ----------------
\noindent\includegraphics[width=\textwidth]{cabecera}\vspace{1cm}%
  \vfill
  % Título proyecto y escudo informática ----------------
  \colorbox{cpardoBox}{%
    \begin{minipage}{.8\textwidth}
      \vspace{.5cm}\Large
      \begin{center}
      \textbf{TFG del Grado en Ingeniería Informática}\vspace{.6cm}\\
      \textbf{\LARGE\@title{}}
      \end{center}
      \vspace{.2cm}
    \end{minipage}

  }%
  \hfill\begin{minipage}{.20\textwidth}
    \includegraphics[width=\textwidth]{escudoInfor}
  \end{minipage}
  \vfill
  % Datos de alumno, curso y tutores ------------------
  \begin{center}%
  {%
    \noindent\LARGE
    Presentado por \@author{}\\ 
    en Universidad de Burgos --- \@date{}\\
    Tutor: \@tutor{}\\
    \@cotutor{}\\
  }%
  \end{center}%
  \null
  \cleardoublepage
  }
\makeatother

\newcommand{\nombre}{Saúl Martín Ibáñez} %%% cambio de comando

% Datos de portada
\title{Estudio de herramientas de realidad aumentada y caso práctico}
\author{\nombre}
\tutor{Pedro Luis Sánchez Ortega}
\cotutor{Álvar Arnaiz González}
\date{\today}

\begin{document}

\maketitle


\newpage\null\thispagestyle{empty}\newpage


%%%%%%%%%%%%%%%%%%%%%%%%%%%%%%%%%%%%%%%%%%%%%%%%%%%%%%%%%%%%%%%%%%%%%%%%%%%%%%%%%%%%%%%%
\thispagestyle{empty}


\noindent\includegraphics[width=\textwidth]{cabecera}\vspace{1cm}

\noindent D. Pedro Luis Sánchez Ortega, profesor del departamento de Ingeniería Electromecánica. Área de Tecnología Electrónica y D. Álvar Arnaiz González del departamento de Ingeniería Informática, área de Lenguajes y Sistemas Informáticos.

\noindent Expone:

\noindent Que el alumno D. \nombre, con DNI 71364596N, ha realizado el Trabajo final de Grado en Ingeniería Informática titulado título de TFG. 

\noindent Y que dicho trabajo ha sido realizado por el alumno bajo la dirección del que suscribe, en virtud de lo cual se autoriza su presentación y defensa.

\begin{center} %\large
En Burgos, {\large \today}
\end{center}

\vfill\vfill\vfill

% Author and supervisor
\begin{minipage}{0.45\textwidth}
\begin{flushleft} %\large
Vº. Bº. del Tutor:\\[2cm]
D. Pedro Luis Sánchez Ortega
\end{flushleft}
\end{minipage}
\hfill
\begin{minipage}{0.45\textwidth}
\begin{flushleft} %\large
Vº. Bº. del co-tutor:\\[2cm]
D. Álvar Arnaiz González
\end{flushleft}
\end{minipage}
\hfill

\vfill

% para casos con solo un tutor comentar lo anterior
% y descomentar lo siguiente
%Vº. Bº. del Tutor:\\[2cm]
%D. nombre tutor


\newpage\null\thispagestyle{empty}\newpage




\frontmatter

% Abstract en castellano
\renewcommand*\abstractname{Resumen}
\begin{abstract}
Desde hace unas décadas, los avances en la informática han hecho que el mundo real conviva con el virtual, llegando incluso a reemplazarlo en algunas tareas o propósitos. Las técnicas que combinan elementos virtuales con el mundo real se ha denomina realidad aumentada.

En este proyecto se va ha realizar un estudio de diferentes herramientas de realidad aumentada actuales. Se compararán sus distintas funcionalidades así como sus pros y contra. Escogiendo cuales de las herramientas estudiadas son las que mas destacan frente al resto. Ademas el estudio se enfocará en cómo estas herramientas de realidad aumentada pueden usarse para la docencia.

Para completar el Trabajo de Fin de Grado se realizará una aplicación orientada a la docencia, con una de las herramientas estudiadas previamente, en la que poner a prueba sus características. La idea de la aplicación es fomentar nuevas formas de aprendizaje mediante el uso de la realidad aumentada, el ejemplo desarrollado trata de ver el proceso de crecimiento de una planta.


\end{abstract}

\renewcommand*\abstractname{Descriptores}
\begin{abstract}
Realidad aumentada, Unity, Vuforia, Aplicación Android, herramientas educativas \ldots
\end{abstract}

\clearpage

% Abstract en inglés
\renewcommand*\abstractname{Abstract}
\begin{abstract}
A \textbf{brief} presentation of the topic addressed in the project.
\end{abstract}

\renewcommand*\abstractname{Keywords}
\begin{abstract}
keywords separated by commas.
\end{abstract}

\clearpage

% Indices
\tableofcontents

\clearpage

\listoffigures

\clearpage

\listoftables
\clearpage

\mainmatter
\capitulo{1}{Introducción}

Las tecnologías de realidad aumentada, consisten en técnicas que consiguen combinar objetos virtuales con objetos del mundo real en directo. 
Estas tecnologías hoy en día están en gran auge en parte por las mejoras de las tecnologías y hardware necesarios. Ya en el 2010 fue seleccionada por la revista \textit{Time} en el top 10 de tendencias tecnológicas, concretamente en el cuarto puesto~\cite{fletcher_2010}. 

Durante la última década ha crecido mucho, hasta el punto en el que actualmente se pueden encontrar aplicadas en múltiples campos de nuestro día a día. Podemos encontrar ejemplos en muchos tipos de aplicaciones de dispositivos móviles y juegos como el \textit{Pokemon Go}. Aplicaciones educativas que ofrecen nuevas técnicas de aprendizaje, como el \textit{MergeCube}, o aplicaciones de navegación como \textit{Google Maps}\footnote{\url{https://support.google.com/maps/thread/12924999?hl=es}}, que añadió un modo de realidad aumentada con el que superponer las indicaciones de la ruta a seguir en la imagen captadas por la cámara del terminal. También hay aplicaciones para uso empresarial como \textit{Vuforia}, que pueden ofrecer apoyo para el mantenimiento y construcción de maquinarias, en la medicina, etc.

Este gran crecimiento es posible gracias al avance en las tecnologías móviles, principalmente, y su accesibilidad a un gran número de personas: los dispositivos móviles tienen la posibilidad de ejecutar tecnologías de realidad aumentada y virtual de forma sencilla y están ``Al alcance de prácticamente todo el mundo``.

\imagen{Memoria/Introducción}{Imágenes de diferentes ejemplos de realidad aumentada~\cite{fotoAR1},~\cite{fotoAR2},~\cite{fotoAR3},~\cite{fotoAR4}.}



Se pueden distinguir dos tipos de realidad aumentada principalmente: los que usan marcadores y los que detectan el entorno. En la detección de marcadores se utilizará una imagen o código QR como eje de referencia para iniciar la realidad aumentada y situar los elementos virtuales en su correcta posición. 
En el segundo tipo la detección del entorno, la cual podrá detectar diferentes planos y paredes reales que podrán usar como referencias, sin la necesidad de una imagen o código QR concreto.


En el proyecto se estudiarán diferentes herramientas de realidad aumentada, comparándolas y, posteriormente, escogeremos una con la que desarrollar un ejemplo en el que poner en práctica las diferentes características y funcionalidades de la herramienta.



\capitulo{2}{Objetivos del proyecto}

Este apartado explica de forma precisa y concisa cuales son los objetivos que se persiguen con la realización del proyecto. Se puede distinguir entre los objetivos marcados por los requisitos del software a construir y los objetivos de carácter técnico que plantea a la hora de llevar a la práctica el proyecto.

El proyecto tiene por objeto proponer entornos de desarrollo de aplicaciones de Realidad Aumentada y Virtual basados en herramientas libres, de código abierto. Se platea conocer las diferentes herramientas disponibles en estos momentos, así como aquellas soluciones más apropiadas para cada caso, móvil, gafas de RV u otros.
\capitulo{3}{Conceptos teóricos}


\section{Realidad aumentada y virtual}
El concepto de realidad aumentada (AR--Augmented Reality) se aplica a las tecnologías que permiten añadir información gráfica sobrepuesta al mundo real en dispositivos visuales: pantallas, gafas, etc.
La realidad virtual (VR--Virtuality Reality) se diferencia de la anterior, en que esta crea un entorno virtual completo en el que el usuario se sumerge.
 
origenes y evolucion...
El termino de 



\section{Herramientas AR}
Se han comparado diferentes herramientas de desarrollo para realidad aumentada. 
Algunas de ellas se tratan de herramientas ya desarrolladas y usadas con fines educativos o comerciales. Otras de las herramientas que se explican son los frameworks/API que muchas de las herramientas anteriores usan.
nuevas online, poder la web.. 

\subsection{ArCore}
ArCore \footnote{\url{https://developers.google.com/ar/discover}} es la plataforma de Google de desarrollo para realidad aumentada. Con el uso de diferentes API permite que los dispositivos detecten su entorno, lo comprendan e interactúen con la información.

Para conseguir integrar el contenido virtual en el mundo real, ArCore utiliza tres técnicas fundamentales.
\begin{itemize}
	\item Motion Tracking (seguimiento del movimiento): permite establecer la posición del móvil en relación con el mundo.
	\item La compresión ambiental: esta permite detectar la posición y tamaño de las superficies del entorno.
	\item Estimación de la luz: que permite estimar las condiciones de luz del ambiente.
\end{itemize}

Su funcionamiento se puede resumir en dos pasos principalmente: rastrear la posición del dispositivo a medida que se mueve y construir su propia comprensión del mundo real.

ArCore utiliza el Motion Tracking para identificar los puntos clave, y rastrea cómo esos puntos se mueven. Combinado el movimiento de dichos puntos con  las lecturas de los sensores de inercia del dispositivo, ArCore calcula la posición y orientación del dispositivo mientras se mueve \cite{google}.

Lamentablemente no es compatible con cualquier dispositivo. En el siguiente enlace se encuentra la lista de los dispositivos compatibles actualmente \url{https://developers.google.com/ar/discover/supported-devices}.

\subsection{ArKit}
ARKit \footnote{\url{https://developer.apple.com/augmented-reality/}} es la herramienta de realidad aumentada de Apple para sistemas iOS. ARKit consigue mostrar contenido virtual de forma natural en el mudo real, incluso pudiendo situarlo detrás o delante del usuario con People Occlusion, pudiendo reconocer hasta 3 rostros al mismo tiempo \cite{apple_inc}.

Funciones destacadas que posee ArKit actualmente.
\begin{itemize}
	\item Oclusión de personas: es capaz de diferenciar a una persona del fondo del escenario, consiguiendo que el contenido virtual  pueda pasar por delante o por detrás de la persona.
	\item Captura de movimiento: Tiene la capacidad de capturar los movimientos de una persona. Distinguiendo diferentes posiciones y movimientos al instante, de forma que es capaz de usarlo como referencias para experiencias de AR.
	\item Cámara frontal y trasera simultánea: Tiene la capacidad de usar ambas cámaras al mismo tiempo, pudiendo así por ejemplo interactuar con el entorno capturado por la cámara trasera, usando unicamente el rostro.
	\item Seguimiento de rostros múltiples: Es capaz de reconocer hasta 3 rostros al mismo tiempo usando la cámara TrueDepth.
	\item Sesiones colaborativas: Tiene la capacidad de crear un mapa mundial, entre múltiples usuarios conectados. Siendo así capaz de crear experiencias de realidad aumentada mas rápido, que pueden servir por ejemplo para juegos multijugador.
\end{itemize}



Una limitación es que unicamente esta disponible para sistemas iOS 11.0 o superior.


\subsection{CoSpaces} CoSpaces\footnote{\url{https://cospaces.io/edu/}} es una plataforma de apoyo educativo. Se puede acceder a ella tanto por web, como desde la aplicación de móvil. Permite a los profesores crear salas para sus alumnos donde tendrían acceso a los diferentes ejemplos a usar para sus clases. También cuenta con una galería donde los usuarios pueden compartir sus ejemplos.
\imagen{cospaces2}{Imagen de mi plataforma de desarrollo de CoSpaces.}
Para la creación de proyectos, se puede hacer desde la web mediante un entorno de desarrollo 3D, donde puedes añadir paredes, objetos y mas tipos de modelos3D, para tu escenario. Con una licencia gratuita, permite guardar dos proyectos como máximo al mismo tiempo. En los proyectos también es posible codificar eventos, acciones mediante una programación de bloques o scripts, aunque hay que destacar, que la programación por scripts y algunas opciones de la programación de bloques únicamente están disponibles con una licencia premium. Para añadir un modelo 3D, que Cospace no ofrezca por defecto, se pueden subir desde un archivo local o desde la búsqueda integrada de CoSpace en Google Poly (biblioteca pública de google de modelos 3D\footnote{\url{https://poly.google.com/}}).
CoSpace también cuenta con la opción de diseñar proyectos para el Mergecube, aunque este plugin se encuentra disponible para la versión premium.

...

\subsection{Metaverse} Metaverse\footnote{\url{https://studio.gometa.io/discover/me}} es otra aplicación enfocada bastante al entorno educativo. En este caso el entorno de desarrollo de los ejemplos es más sencillo. En la propia web, tendremos una estructura similar un modelo de diagrama de flujo. En cada paso se puede añadir un objeto 3D o una imagen que es la que estará flotando cuando estemos usando la realidad aumentada, posteriormente se pueden añadir botones y menús, para que te lleven a otra pantalla que tenga otro objeto asignado. Así, por ejemplo, se pueden hacer programas en los que se creen pequeños juegos de preguntas y, dependiendo de las respuestas, te llevan a diferentes pantallas.
También da la posibilidad de importar tus propios modelos e imágenes.
También incluye la posibilidad de reconocer expresiones faciales para poder usarlo como disparadores.
También puede usar el ArCore y ArKit para permitir que los objetos puedan asentarse en una posición y poder girar alrededor suyo.

\subsection{Zapworks} ZapWorks\footnote{\url{https://zap.works/}} se trata de una herramienta de desarrollo para AR. En general esta compuesto por dos herramientas de creación, ZapWorks Designer y ZapWorks Studio.
\begin{itemize}
	\item ZapWorks Studio: se trata de la aplicación de desarrollo para escritorio. La aplicación permite desarrollar los proyectos AR, con modelos 3D, animaciones, interacciones, etc. Soporta los formatos  GLTF, FBX, OBJ, POD de modelos 3D.
	\imagen{zapworkStudio}{Imagen perteneciente a la aplicación de ZapWorks Studio\cite{ZapStudio}.}
	\item ZapWorks Designer: Se trata de una aplicación web, desde la que se puede de una forma sencilla, asociar a un marcador, imágenes, textos, vídeos, links. Pero no es posible trabajar con modelos 3D.
\end{itemize}	
La herramienta posee los siguientes tipos de seguimiento de realidad aumentada.
\begin{itemize}
	\item Word Tracking: Capacidad de reconocer el entorno para poder colocar objetos de una manera mas realista sobre el terreno, sin la ayuda de marcadores.
	\item Face Tracking: Capacidad de reconocer y seguir rostros, para por ejemplo colocar, en una cara el contenido AR y que aunque se muevan el contenido se mueva con la cara.
	\item Image Tracking: Que una imagen pueda ser el marcador, y pueda seguirla en movimiento.
	\item Zapbox Tracking: \colorbox{red}{COMPLETAR}
	 
\end{itemize}
Para acceder al contenido AR almacenado en el servidor de ZapWorks, utiliza los zapcodes, se tratan una imagen codificada para que pueda descargar un proyecto AR. Una vez que la aplicación reconoce el zapcode, descarga el proyecto y este se ejecutará.
\imagenPeque{zapcode}{Ejemplo de un zapcode \cite{ZapStudio}.}


\subsection{Kudan} Kudan\footnote{\url{https://www.xlsoft.com/en/products/kudan/index.html}} 
se trata de un SDK, con la capacidad de soportar AR tanto con marcadores y como sin ellos. En principio Kudan no tiene un limite de marcadores que pueda detectar al mismo tiempo, pero se puede incluir un limite para que la aplicación tenga un rendimiento adecuado. 
Posee soporte para APIs nativas, como ObjetiveC para iOS, Java para Android y Unity \cite{kudan_developer_hub}.

El motor de Kudan está escrito en C++ y optimizado con programación en ensamblador, dando un mayor rendimiento y estabilidad, con el menor impacto de memoria.
Usa la tecnología SLAM (Simultaneous Localization and Mapping) para el reconocimiento del escenario.

\imagenPeque{marta_2}{Imagen perteneciente a la aplicación M.A.R.T.A, que funciona con Kudan}

Kudan soporta los siguientes formatos 3D: FBX, OBJ y COLLADA.

\subsection{Vuforia}
Vuforia\footnote{\url{https://developer.vuforia.com/}} es un kit de desarrollo ( SDK) para aplicaciones de AR, para las plataformas de Android, iOS, Windows y HoloLens. Fue creado por Qualcomm Connected Experience en 2010, y en 2015 PTC inc lo compró \cite{simonetti2013vuforia}.

Vuforia proporciona una API en C++, Java, Objective-C++ y los lenguajes de .NET mediante Unity.

Vuforia soporta diferentes tipos marcadores, estos pueden ser 2D o 3D, también soporta múltiples marcadores simultáneamente, y reconocimiento del terreno sin marcadores.
\imagen{vuforia01}{Posibles marcadores de Vuforia}

Los marcadores se pueden crear desde la página web de developer vuforia, para ello deben escoger qué tipo de marcadores se quieren crear y adjuntarle las imágenes que formarán dicho marcador. La página dará una calificación de 5 estrellas según la calidad de la imagen para ser un marcador. Una vez completado, da la opción de descargarlos el un archivo configurable de Unity que incluirá los marcadores en nuestro proyecto.



Cuenta con una licencia gratuita por defecto, con la limitación de poder tener un máximo de 100 vumarks y 1000 cloud targets.

Actualmente Vuforia esta disponible para las versiones indicadas en la siguiente tabla~\cite{vuforia_supported_versions}.
\begin{table}[]
	\begin{tabular}{|l|l|l|l|l|l|}
		\hline
		\multicolumn{2}{|l|}{Device OS}                       & \multicolumn{2}{l|}{Herramientas de Desarrollo}                                             & \multicolumn{2}{l|}{Fusion Provider}                \\ \hline
		&                          & NDK                                            & r20+                                       &                                &                    \\ \cline{3-4}
		&                          & \cellcolor[HTML]{EFEFEF}Gradle                 & \cellcolor[HTML]{EFEFEF}5.1.1+             &                                &                    \\ \cline{3-4}
		&                          & Android SDK Build Tools                        & 28.0.3                                     &                                &                    \\ \cline{3-4}
		&                          & \cellcolor[HTML]{EFEFEF}Android Studio         & \cellcolor[HTML]{EFEFEF}3.4.x              &                                &                    \\ \cline{3-4}
		\multirow{-5}{*}{Android}  & \multirow{-5}{*}{5.1.1+} & Unity Editor                                   & 2019.2.0+                                  & \multirow{-5}{*}{ArCore 1.10+} & \multirow{-5}{*}{} \\ \hline
		&                          & \cellcolor[HTML]{EFEFEF}XCode                  & \cellcolor[HTML]{EFEFEF}10.1+              & \multicolumn{2}{l|}{}                               \\ \cline{3-4}
		\multirow{-2}{*}{iOS}      & \multirow{-2}{*}{11+}    & Unity Editor                                   & 2019.2+                                    & \multicolumn{2}{l|}{\multirow{-2}{*}{ArKit}}        \\ \hline
		&                          & \cellcolor[HTML]{EFEFEF}Visual Studio          & \cellcolor[HTML]{EFEFEF}2017 v 15.9+ & \multicolumn{2}{l|}{}                               \\ \cline{3-4}
		&                          & Unity Editor                                   & 2019.2.0+                                  & \multicolumn{2}{l|}{}                               \\ \cline{3-4}
		\multirow{-3}{*}{Windows}  & \multirow{-3}{*}{10}     & \cellcolor[HTML]{EFEFEF}Unity Editor(Hololens) & \cellcolor[HTML]{EFEFEF}2018.4.11          & \multicolumn{2}{l|}{\multirow{-3}{*}{}}             \\ \hline
		&                          & Lumin SDK                                      & 0.22.0                                     & \multicolumn{2}{l|}{}                               \\ \cline{3-4}
		\multirow{-2}{*}{Lumin Os} & \multirow{-2}{*}{10}     & \cellcolor[HTML]{EFEFEF}Lumin OS               & \cellcolor[HTML]{EFEFEF}0.97+              & \multicolumn{2}{l|}{\multirow{-2}{*}{}}             \\ \hline
	\end{tabular}
\end{table}


\subsection{Wikitude}
Wikitude\footnote{\url{https://www.wikitude.com/}} se trata de un SDK de realidad aumentada, ofrece un amplio repertorio de características\cite{wikitude}: 
\begin{itemize}
	\item Reconocimiento de imágenes.
	\item Reconocimiento de objetos.
	\item Reconocimiento de múltiples marcadores: Permite reconocer múltiples marcadores de forma simultanea.
	\item Instant Tracking: reconocimiento sin la necesidad de marcadores, permitiendo superponer elementos virtuales en superficies detectadas con la tecnologia Slam Instant Tracking.
	
	\item Seguimiento extendido: Permite que un elemento virtual que se ha superpuesto en un punto o marcador determinado, persista aunque dicho punto se salga del campo de visión de la cámara.
	
	\item Geo AR: Permite agregar contenido de realidad aumentada basándose en la ubicación, valiéndose del GPS y demás sensores para determinar la ubicación.
	\item Cloud Recognition: Ofrece la posibilidad de guardar en linea los datos de imágenes/marcadores para las aplicaciones.
\end{itemize}

Wikitude esta disponible para Android, iOS, Windows y Smart Glasses.
También posee un amplio soporte para frameworks de desarrollo : Andoid, iOS, Windows, Unity, Cordova, Xamarin, Flutter, Titanium. 
Los sistemas deberán cumplir los requisitos indicados en el siguiente enlace:\url{https://www.wikitude.com/documentation/latest/android/supporteddevices.html#supported-devices}.

\subsection{ArUco}


\subsection{OpenCV}
OpenCV\footnote{\url{https://opencv.org/}} es una librería de visión artificial y machine learning, de código abierto. Tiene una licencia BSD por lo que facilita su uso y modificación para las empresas. Cuenta con más de 2500 algoritmos optimizados. Estos pueden emplearse para detectar y reconocer rostros, identificación de objetos, clasificar acciones humanas en vídeos, seguir los movimientos de la cámara, rastrear objetos en movimiento, extraer modelos 3D de objetos, producir nubes de puntos 3D de cámaras estéreo, seguir el movimiento de los ojos, reconocer paisajes y establecer marcadores para usarlos en realidad aumentada.

Esta disponible en interfaces para C++,Python, Java y MATLAB y es compatible con Windows, Linux, Android y Mac OS.


\subsection{8thWall}
8thWall\footnote{\url{https://www.8thwall.com/}}
es un entorno de desarrollo de realidad aumentada, que se caracteriza principalmente por el WebAR, que da la posibilidad de ejecutar las imágenes AR a través del propio navegador web del dispositivo, sin tener que instalarse ninguna aplicación extra en el mismo\cite{8thwall_products}.

Posee las siguientes características:
\begin{itemize}
	\item World Tracking: (COMPLETAR)
	\item Image Targets: 8thWall Web puede detectar y rastrear imágenes y usarlas como marcadores. Cada aplicación puede tener un máximo de 1000 marcadores.
	\item Modular framework: (COMPLETAR)
\end{itemize}


Tiene soporte con los frameworks 3D A-Frame, three.js, babylon.js, Amazon Sumerian y PlayCanvas.

Requiere en iOS, iOS11 o superior, y el navegador Safari. En Android los navegadores soportados son; Chrome, Chromium,Firefox, Android WebViews.

\section{Técnicas AR}
	\subsection{Detección de marcadores}
	La detección de marcadores se trata de un técnica de realidad aumentada en el que los sistemas utilizan como referencia espacial, para situar elementos que queremos superponer, una figura o imagen concreta, que ha sido especificada previamente. Los marcadores mas comunes y simples son los códigos QR, pues se tratan de imágenes diseñadas para que puedan ser reconocidas fácilmente por las cámaras.
	\imagen{ejemploMarcadorCasa}{Ejemplo un marcador sencillo  \cite{AR_company}.}
	Con la mejora de las cámaras, también es posible utilizar como marcadores imágenes mas complejas, como por ejemplo un folleto publicitario o la foto de un planeta.
	También es posible que el marcador se trate de un objeto físico, como por ejemplo una lámpara, o un coche.
	\imagen{ejemploMarcadorTigre}{Ejemplo donde una carta de un tigre es un marcador \cite{AR_company}.}
	Aunque gracias al avance de la calidad de las cámaras y del software de reconocimiento, es posible utilizar marcadores mas complejos, esto también implicarán que será mas probable que falle, y por consecuencia se desencuandren las imágenes insertadas por AR, que los tiempos de detección sean mas largos etc.
	
	Se pueden clasificar los marcadores en \cite{linowes_babilinski_2017}:
	\begin{itemize}
		\item \subsubsection{Marcadores}El marcador mas básico es uno con un borde ancho. Su ventaja es que minimizan los costes de procesamiento y las posibilidades de error al reconocerlo.
		\imagenPeque{market}{Ejemplo de un marcador\cite{linowes_babilinski_2017}.}
		\item \subsubsection{Marcadores Codificados}Se trata de marcadores que poseen códigos de barra 2D, dentro del área de bordes patrones. 
		\imagen{market2}{Ejemplo de marcadores codificados\cite{linowes_babilinski_2017}.}
		\item \subsubsection{Imágenes}El poder usar imágenes evita tener que crear marcadores personalizados. El reconocimiento de imágenes entra en la categoría de seguimiento de características naturales. Su calidad como marcador, dependerá de si tiene un borde que defina claramente esa imagen, del contraste de sus colores o la complejidad de su composición.
		\item \subsubsection{Multiples Marcadores}
		Es el caso de usar mas de un marcador al mismo tiempo por la misma cámara. Permitiendo así tener distintos objetos en la misma escena simultáneamente. También es posible crearlos formado objetos geométricos, un ejemplo sería el Mergecube(Inluir referencia merge).
		\item \subsubsection{Reconocimiento de texto}
		Algunos SDK de realidad aumentada cuentan con la capacidad de distinguir textos. La precisión estará influenciada dependiendo del tipo de fuente de las letra. 
		\item \subsubsection{Formas Simples}
		Algunas aplicaciones tienen la capacidad de distinguir diferentes formas simples, como son cubos, cilindros. También pueden distinguir las medidas de las formas, pudiendo así diferenciar si un cubo es mas grande que otro.
		\item \subsubsection{Reconocimiento de objetos}
		Se trata del siguiente nivel, la capacidad de distinguir objetos. No solo reconocerá e identifica la forma y tamaño del objeto, también detalles característicos que le puedan diferenciar de otro. Vuforia por ejemplo cuenta con una funcionalidad llamada Vuforia Object Scanner, que permite escanear un objeto para una copia digital, que se podría usar como marcador.
		\imagen{vuforia_recoObjeto}{Vuforia Object Scanner \cite{vuforia_scanner}.}		
	\end{itemize}
	\subsubsection{Mergecube}
	El Mergecube\footnote{\url{https://mergeedu.com/}} se trata de un cubo diseñado por Merge, el cual tiene grabado unos dibujos por sus 6 caras, los dibujos actúan de marcadores. Es posible mover el cubo al mismo tiempo que se usa, de esa manera se consigue una experiencia mas interactiva. Es por eso que el mergecube es ideal para aplicarlo en experiencias educativas.
	\imagen{mergecube}{Mergecube  \cite{barker_2019}.} 
	\subsection{Detección del entorno}
	La detección del entorno trata de poder reconocer el entorno que la cámara capta y ser capaz de, sin la ayuda de un marcador, poder ubicar los modelos en el escenario.
	Para esto hay diferentes técnicas, en resumen estas técnicas lo que hacen es detectar los cambios de profundidad de la imagen, y con esos datos evaluar qué objetos están mas o menos cerca, cuáles son planos, si se trata de un objeto que no tiene profundidad(una pared)...
	\subsubsection{SLAM}
	Simultaneous Localization and Mapping, se trata de una técnica para estimar los movimientos de los sensores y reconstruir la estructura en un entorno desconocido. Cuando esta técnica se aplica en cámaras, se denomina visual SLAM(vSLAM), ya que se basara en la información visual recibida.
	La mayoría de herramientas de realidad aumentada, que tienen la capacidad de detectar el entorno sin necesidad de marcadores, se basan en esta técnica.
	
	Los algoritmos de vSLAM, no solo son usados unicamente para la AR, también pueden ser aplicados para robótica, o los coches autónomos no tripulados.
	
	
\section{Secciones}

Las secciones se incluyen con el comando section.

\subsection{Subsecciones}

Además de secciones tenemos subsecciones.

\subsubsection{Subsubsecciones}

Y subsecciones. 


\section{Referencias}

Las referencias se incluyen en el texto usando cite \cite{wiki:latex}. Para citar webs, artículos o libros \cite{koza92}.


\section{Imágenes}

Se pueden incluir imágenes con los comandos standard de \LaTeX, pero esta plantilla dispone de comandos propios como por ejemplo el siguiente:

\imagen{escudoInfor}{Autómata para una expresión vacía}



\section{Listas de items}

Existen tres posibilidades:

\begin{itemize}
	\item primer item.
	\item segundo item.
\end{itemize}

\begin{enumerate}
	\item primer item.
	\item segundo item.
\end{enumerate}

\begin{description}
	\item[Primer item] más información sobre el primer item.
	\item[Segundo item] más información sobre el segundo item.
\end{description}
	
\begin{itemize}
\item 
\end{itemize}

\section{Tablas}

Igualmente se pueden usar los comandos específicos de \LaTeX o bien usar alguno de los comandos de la plantilla.

\tablaSmall{Herramientas y tecnologías utilizadas en cada parte del proyecto}{l c c c c}{herramientasportipodeuso}
{ \multicolumn{1}{l}{Herramientas} & App AngularJS & API REST & BD & Memoria \\}{ 
HTML5 & X & & &\\
CSS3 & X & & &\\
BOOTSTRAP & X & & &\\
JavaScript & X & & &\\
AngularJS & X & & &\\
Bower & X & & &\\
PHP & & X & &\\
Karma + Jasmine & X & & &\\
Slim framework & & X & &\\
Idiorm & & X & &\\
Composer & & X & &\\
JSON & X & X & &\\
PhpStorm & X & X & &\\
MySQL & & & X &\\
PhpMyAdmin & & & X &\\
Git + BitBucket & X & X & X & X\\
Mik\TeX{} & & & & X\\
\TeX{}Maker & & & & X\\
Astah & & & & X\\
Balsamiq Mockups & X & & &\\
VersionOne & X & X & X & X\\
} 

\capitulo{4}{Técnicas y herramientas}

\section{Herramientas}
\subsection{Herramientas AR}

Se han comparado diferentes herramientas de desarrollo para realidad aumentada. 
Algunas de ellas se tratan de herramientas ya desarrolladas y usadas con fines educativos o comerciales. Otras de las herramientas que se explican son los frameworks/API que muchas de las herramientas anteriores usan.


\subsubsection{ArCore}

ArCore \footnote{\url{https://developers.google.com/ar/discover}} es la plataforma de Google de desarrollo para realidad aumentada. Con el uso de diferentes API permite que los dispositivos detecten su entorno, lo comprendan e interactúen con la información.

Para conseguir integrar el contenido virtual en el mundo real, ArCore utiliza tres técnicas fundamentales:
\begin{itemize}
	\item Motion Tracking (seguimiento del movimiento): permite establecer la posición del móvil en relación con el mundo.
	\item La compresión ambiental: esta permite detectar la posición y tamaño de las superficies del entorno.
	\item Estimación de la luz: que permite calcular las condiciones de luz del ambiente.
\end{itemize}

Su funcionamiento se puede resumir en dos pasos principalmente: rastrear la posición del dispositivo a medida que se mueve y construir su propia comprensión del mundo real.

ArCore utiliza el Motion Tracking para identificar los puntos clave, y rastrea cómo esos puntos se mueven. Combinado el movimiento de dichos puntos con  las lecturas de los sensores de inercia del dispositivo, ArCore calcula la posición y orientación del dispositivo mientras se mueve~\cite{google}.

Lamentablemente no es compatible con cualquier dispositivo. En el siguiente enlace se encuentra la lista de los dispositivos compatibles actualmente \url{https://developers.google.com/ar/discover/supported-devices}.

\subsubsection{ArKit}

ARKit \footnote{\url{https://developer.apple.com/augmented-reality/}} es la herramienta de realidad aumentada de Apple para sistemas iOS. ARKit consigue mostrar contenido virtual de forma natural en el mudo real, incluso pudiendo situarlo detrás o delante del usuario con People Occlusion, pudiendo reconocer hasta 3 rostros al mismo tiempo~\cite{apple_inc}.

Funciones destacadas que posee ArKit actualmente.

\begin{itemize}
	\item Oclusión de personas: es capaz de diferenciar a una persona del fondo del escenario, consiguiendo que el contenido virtual  pueda pasar por delante o por detrás de la persona.
	\item Captura de movimiento: Tiene la capacidad de capturar los movimientos de una persona. Distinguiendo diferentes posiciones y movimientos al instante, de forma que es capaz de usarlo como referencias para experiencias de AR.
	\item Cámara frontal y trasera simultánea: Tiene la capacidad de usar ambas cámaras al mismo tiempo, pudiendo así por ejemplo interactuar con el entorno capturado por la cámara trasera, usando únicamente el rostro.
	\item Seguimiento de rostros múltiples: Es capaz de reconocer hasta 3 rostros al mismo tiempo usando la cámara TrueDepth.
	\item Sesiones colaborativas: Tiene la capacidad de crear un mapa mundial, entre múltiples usuarios conectados. Siendo así capaz de crear experiencias de realidad aumentada mas rápido, que pueden servir por ejemplo para juegos multijugador.
\end{itemize}



Una limitación es que únicamente esta disponible para sistemas iOS 11.0 o superior.


\subsubsection{CoSpaces}\label{sub:Def_cospace} 

CoSpaces\footnote{\url{https://cospaces.io/edu/}} es una plataforma de apoyo educativo. Se puede acceder a ella tanto por web, como desde la aplicación de móvil. Permite a los profesores crear salas para sus alumnos donde tendrían acceso a los diferentes ejemplos a usar para sus clases. También cuenta con una galería donde los usuarios pueden compartir sus ejemplos.

\imagen{cospaces2}{Imagen de mi plataforma de desarrollo de CoSpaces. ( Captura de pantalla)}

Para la creación de proyectos, se puede hacer desde la web mediante un entorno de desarrollo 3D, donde puedes añadir paredes, objetos y mas tipos de modelos3D, para tu escenario. Con una licencia gratuita, permite guardar dos proyectos como máximo al mismo tiempo, en la imagen~\ref{fig:cospaces2} se puede observar el entorno descrito. En los proyectos también es posible codificar eventos, acciones mediante una programación de bloques o scripts, aunque hay que destacar, que la programación por scripts y algunas opciones de la programación de bloques únicamente están disponibles con una licencia premium. Para añadir un modelo 3D, que Cospace no ofrezca por defecto, se pueden subir desde un archivo local o desde la búsqueda integrada de CoSpace en Google Poly (biblioteca pública de google de modelos 3D\footnote{\url{https://poly.google.com/}}).
CoSpace también cuenta con la opción de diseñar proyectos para el Mergecube, aunque este plugin se encuentra disponible para la versión premium.



\subsubsection{Metaverse} 

Metaverse\footnote{\url{https://studio.gometa.io/discover/me}} es otra aplicación enfocada bastante al entorno educativo. En este caso el entorno de desarrollo de los ejemplos es más sencillo. En la propia web, tendremos una estructura similar a un modelo de diagrama de flujo. En cada paso se puede añadir un objeto 3D o una imagen que es la que estará flotando cuando estemos usando la realidad aumentada, posteriormente se pueden añadir botones y menús, para que te lleven a otra pantalla que tenga otro objeto asignado. Así, por ejemplo, se pueden hacer programas en los que se creen pequeños juegos de preguntas y, dependiendo de las respuestas, te llevan a diferentes pantallas.
Además da la posibilidad de importar tus propios modelos e imágenes.
También incluye la posibilidad de reconocer expresiones faciales para poder usarlo como disparadores.
Otra característica es que puede usar el ArCore y ArKit para permitir que los objetos puedan asentarse en una posición y poder girar alrededor suyo.

\subsubsection{Zapworks} ZapWorks\footnote{\url{https://zap.works/}} se trata de una herramienta de desarrollo para AR. En general esta compuesto por dos herramientas de creación, ZapWorks Designer y ZapWorks Studio.

\begin{itemize}
	\item ZapWorks Studio: se trata de la aplicación de desarrollo para escritorio. La aplicación permite desarrollar los proyectos AR, con modelos 3D, animaciones, interacciones, etc. Soporta los formatos  GLTF, FBX, OBJ, POD de modelos 3D. En la imagen~\ref{fig:zapworkStudio} se puede observar el entorno.
	\item ZapWorks Designer: Se trata de una aplicación web, desde la que se puede de una forma sencilla, asociar a un marcador, imágenes, textos, vídeos, links. Pero no es posible trabajar con modelos 3D.
\end{itemize}	

\imagen{zapworkStudio}{Imagen perteneciente a la aplicación de ZapWorks Studio~\cite{ZapStudio}.}

La herramienta posee los siguientes tipos de seguimiento de realidad aumentada.

\begin{itemize}
	\item Word Tracking: Capacidad de reconocer el entorno para poder colocar objetos de una manera mas realista sobre el terreno, sin la ayuda de marcadores.
	\item Face Tracking: Capacidad de reconocer y seguir rostros, para por ejemplo colocar, en una cara el contenido AR y que aunque se muevan el contenido se mueva con la cara.
	\item Image Tracking: Que una imagen pueda ser el marcador, y pueda seguirla en movimiento.
	\item Zapbox Tracking: Se trata de unos marcadores específicamente creados para que se puedan manipular~\ref{fig:Memoria/zapbox}.
\end{itemize}

\imagen{Memoria/zapbox}{Imagen perteneciente a un kit de zapbox~\cite{ZapStudio}.}

Para acceder al contenido AR almacenado en el servidor de ZapWorks, utiliza los zapcodes~\ref{fig:zapcode}, se tratan una imagen codificada para que pueda descargar un proyecto AR. Una vez que la aplicación reconoce el zapcode, descarga el proyecto y este se ejecutará.

\imagenPeque{zapcode}{Ejemplo de un zapcode~\cite{ZapStudio}.}


\subsubsection{Kudan} 

Kudan\footnote{\url{https://www.xlsoft.com/en/products/kudan/index.html}} 
se trata de un SDK, con la capacidad de soportar AR tanto con marcadores y como sin ellos. En principio Kudan no tiene un limite de marcadores que pueda detectar al mismo tiempo, pero se puede incluir un limite para que la aplicación tenga un rendimiento adecuado. 
Posee soporte para APIs nativas, como ObjetiveC para iOS, Java para Android y Unity~\cite{kudan_developer_hub}.

El motor de Kudan está escrito en C++ y optimizado con programación en ensamblador, dando un mayor rendimiento y estabilidad, con el menor impacto de memoria.
Usa la tecnología SLAM (Simultaneous Localization and Mapping) para el reconocimiento del escenario.

Kudan soporta los siguientes formatos 3D: FBX, OBJ y COLLADA.

En la imagen \ref{fig:marta_2} podemos ver un ejemplo de la aplicación M.A.R.T.A\footnote{\url{https://play.google.com/store/apps/details?id=com.apophistechlabs.marta}} que utiliza Kudan para detectar una superficie en la que colocar el modelo 3D.

\imagenPeque{marta_2}{Imagen perteneciente a la aplicación M.A.R.T.A~\cite{marta_app}.}



\subsubsection{Vuforia}\label{sub:Def_Vuforia}

Vuforia\footnote{\url{https://developer.vuforia.com/}} es un kit de desarrollo ( SDK) para aplicaciones de AR, para las plataformas de Android, iOS, Windows, Unity y HoloLens. Fue creado por Qualcomm Connected Experience en 2010, y en 2015 PTC inc lo compró~\cite{simonetti2013vuforia}.

Vuforia proporciona una API en C++, Java, Objective-C++ y los lenguajes de .NET mediante Unity.

Vuforia soporta diferentes tipos marcadores~\ref{fig:vuforia01}, estos pueden ser 2D o 3D, también soporta múltiples marcadores simultáneamente, y reconocimiento del terreno sin marcadores.

\imagen{vuforia01}{Posibles marcadores de Vuforia}

Los marcadores se pueden crear desde la página web de developer vuforia, para ello deben escoger qué tipo de marcadores se quieren crear y adjuntarle las imágenes que formarán dicho marcador. La página dará una calificación de 5 estrellas según la calidad de la imagen para ser un marcador. Una vez completado, da la opción de descargarlos el un archivo configurable de Unity que incluirá los marcadores en nuestro proyecto.



Cuenta con una licencia gratuita por defecto, con la limitación de poder tener un máximo de 100 vumarks y 1000 cloud targets.

Actualmente Vuforia esta disponible para las versiones indicadas en la siguiente tabla~\cite{vuforia_supported_versions}.

REVISAR TABLA   ELIMINAR SI REALMENTE NO ES RELEVANTE
\begin{table}[]
	\begin{tabular}{|l|l|l|l|l|l|}
		\hline
		\multicolumn{2}{|l|}{Device OS}                       & \multicolumn{2}{l|}{Herramientas de Desarrollo}                                             & \multicolumn{2}{l|}{Fusion Provider}                \\ \hline
		&                          & NDK                                            & r20+                                       &                                &                    \\ \cline{3-4}
		&                          & \cellcolor[HTML]{EFEFEF}Gradle                 & \cellcolor[HTML]{EFEFEF}5.1.1+             &                                &                    \\ \cline{3-4}
		&                          & Android SDK Build Tools                        & 28.0.3                                     &                                &                    \\ \cline{3-4}
		&                          & \cellcolor[HTML]{EFEFEF}Android Studio         & \cellcolor[HTML]{EFEFEF}3.4.x              &                                &                    \\ \cline{3-4}
		\multirow{-5}{*}{Android}  & \multirow{-5}{*}{5.1.1+} & Unity Editor                                   & 2019.2.0+                                  & \multirow{-5}{*}{ArCore 1.10+} & \multirow{-5}{*}{} \\ \hline
		&                          & \cellcolor[HTML]{EFEFEF}XCode                  & \cellcolor[HTML]{EFEFEF}10.1+              & \multicolumn{2}{l|}{}                               \\ \cline{3-4}
		\multirow{-2}{*}{iOS}      & \multirow{-2}{*}{11+}    & Unity Editor                                   & 2019.2+                                    & \multicolumn{2}{l|}{\multirow{-2}{*}{ArKit}}        \\ \hline
		&                          & \cellcolor[HTML]{EFEFEF}Visual Studio          & \cellcolor[HTML]{EFEFEF}2017 v 15.9+ & \multicolumn{2}{l|}{}                               \\ \cline{3-4}
		&                          & Unity Editor                                   & 2019.2.0+                                  & \multicolumn{2}{l|}{}                               \\ \cline{3-4}
		\multirow{-3}{*}{Windows}  & \multirow{-3}{*}{10}     & \cellcolor[HTML]{EFEFEF}Unity Editor(Hololens) & \cellcolor[HTML]{EFEFEF}2018.4.11          & \multicolumn{2}{l|}{\multirow{-3}{*}{}}             \\ \hline
		&                          & Lumin SDK                                      & 0.22.0                                     & \multicolumn{2}{l|}{}                               \\ \cline{3-4}
		\multirow{-2}{*}{Lumin Os} & \multirow{-2}{*}{10}     & \cellcolor[HTML]{EFEFEF}Lumin OS               & \cellcolor[HTML]{EFEFEF}0.97+              & \multicolumn{2}{l|}{\multirow{-2}{*}{}}             \\ \hline
	\end{tabular}
\end{table}


\subsubsection{Wikitude}

Wikitude\footnote{\url{https://www.wikitude.com/}} se trata de un SDK de realidad aumentada, ofrece un amplio repertorio de características~\cite{wikitude}: 

\begin{itemize}
	\item Reconocimiento de imágenes.
	\item Reconocimiento de objetos.
	\item Reconocimiento de múltiples marcadores: Permite reconocer múltiples marcadores de forma simultánea.
	\item Instant Tracking: reconocimiento sin la necesidad de marcadores, permitiendo superponer elementos virtuales en superficies detectadas con la tecnologia Slam Instant Tracking.	
	\item Seguimiento extendido: Permite que un elemento virtual que se ha superpuesto en un punto o marcador determinado, persista aunque dicho punto se salga del campo de visión de la cámara.	
	\item Geo AR: Permite agregar contenido de realidad aumentada basándose en la ubicación, valiéndose del GPS y demás sensores para determinar la ubicación.
	\item Cloud Recognition: Ofrece la posibilidad de guardar en línea los datos de imágenes/marcadores para las aplicaciones.	
\end{itemize}

Wikitude esta disponible para Android, iOS, Windows y Smart Glasses.
También posee un amplio soporte para frameworks de desarrollo : Andoid, iOS, Windows, Unity, Cordova, Xamarin, Flutter, Titanium. 
Los sistemas deberán cumplir los requisitos indicados en el siguiente enlace: \url{https://www.wikitude.com/documentation/latest/android/supporteddevices.html#supported-devices}.



\subsubsection{OpenCV}

OpenCV\footnote{\url{https://opencv.org/}} es una librería de visión artificial y machine learning, de código abierto. Tiene una licencia BSD por lo que facilita su uso y modificación para las empresas. Cuenta con más de 2500 algoritmos optimizados. Estos pueden emplearse para detectar y reconocer rostros, identificación de objetos, clasificar acciones humanas en vídeos, seguir los movimientos de la cámara, rastrear objetos en movimiento, extraer modelos 3D de objetos, producir nubes de puntos 3D de cámaras estéreo, seguir el movimiento de los ojos, reconocer paisajes y establecer marcadores para usarlos en realidad aumentada.

Esta disponible en interfaces para C++, Python, Java y MATLAB y es compatible con Windows, Linux, Android y Mac OS.


\subsubsection{8thWall}

8thWall\footnote{\url{https://www.8thwall.com/}}
es un entorno de desarrollo de realidad aumentada, que se caracteriza principalmente por el WebAR, que da la posibilidad de ejecutar las imágenes AR a través del propio navegador web del dispositivo, sin tener que instalarse ninguna aplicación extra en el mismo~\cite{8thwall_products}.

Posee las siguientes características:

\begin{itemize}
	\item World Tracking: Con la tecnología SLAM, es capaz de reconocer superficies planas instantáneamente, ademas de estimaciones de iluminación.
	\item Image Targets: 8thWall Web puede detectar y rastrear imágenes y usarlas como marcadores. Cada aplicación puede tener un máximo de 1000 marcadores.
	\item Modular framework: El framework de la aplicación esta diseñado para integrar tecnologías de reconocimiento como el seguimiento de rostros, la oclusión de personas y otras tecnologías de machine learning.
\end{itemize}


Tiene soporte con los frameworks 3D A-Frame, three.js, babylon.js, Amazon Sumerian y PlayCanvas.

Requiere en iOS, iOS11 o superior, y el navegador Safari. En Android los navegadores soportados son; Chrome, Chromium,Firefox, Android WebViews.

\subsubsection{Resumen y comparativa de herramientas}
En la tabla \ref{tabla:comparacionHerramientas} se puede observar una comparativa entre las principales funcionalidades de las herramientas. Se indica si es MultiTarget (MT), si posee realidad sin marcadores (SM), las plataformas para las que está disponible, el tipo de licencia, libre (L) o comercial (C), el precio y tipo de requisitos.

\label{tablaComparacion}\tablaSmall{Comparativa herramientas AR}{l c c c c r c}{comparacionHerramientas}
{ \multicolumn{1}{l}{} & MT & SM & Plataforma  & Licencia & Precio & Requisitos\\}{ 
	ArCore & X&X&iOS/Android& L & 0\euro & Altos\\ 
	ArKit &X&X&iOS &L & 0\euro &-\\
	CoSpaces&X & &iOS/Android  &L y C & 75\euro/año&Medios\\
	Metaverse &X&X&iOS/Android &L y C & ?\euro&Medios \\
	Zapworks  &X&&iOS/Android &L y C & 55\euro/mes& Medios\\
	Kudan   &X&X&iOS/Android &L y C & ?\euro& Medios\\
	Vuforia  &X&X& iOS/Android&L y C & 42\euro/mes& Medios\\
	WikiTude &X&X&iOS/Android &L y C & 2490\euro & Medios\\
	OpenCV  &X& & iOS/Android&L & 0\euro& Bajos\\
	8thWall &X&X& iOS/Android&L y C & 99mes\euro& Medios\\
}

\subsection{Herramientas de documentación}
\subsubsection{\LaTeX} 
Para documentar este proyecto se ha utilizado \LaTeX. Se trata de un sistema de composición de textos, con el objetivo de la creación de documentos con una alta calidad tipográfica. Se ha usado TeXstudio y miktex para compilar y editar \LaTeX.

A diferenciara de otros editores clásicos de textos para documentación, \LaTeX tiene una curva de aprendizaje más elevada, pero al final sus ventajas a la hora de documentar, especialmente este tipo de proyectos, le convierte una gran opción.


\subsubsection{Draw.io}
Draw.io\footnote{\url{https://www.diagrams.net/}} se trata de una herramienta online, para el modelado de diagramas UML. Una de sus grandes ventajas es la posibilidad de guardar y sincronizar los diagramas en servicios como Google Drive, OneDrive, Dropbox, GitHub, GitLab y Trello. Funciona de manera gráfica pudiendo arrastrar, conectar y alinear los componentes de los diagramas de forma sencilla.

Actualmente Draw.io se encuentra en proceso de transición a diagrams.net, este cambio no implica ningúna diferencia en el funcionamiento de la aplicación~\cite{diagramsNet}.



\subsection{Herramientas de gestión}

\subsubsection{GitHub}\label{github}
Github\footnote{\url{https://github.com}} se trata de una plataforma en la que es posible alojar proyectos usando un control de versiones Git.

Se ha utilizado esta herramienta para alojar el repositorio del proyecto desarrollado, ademas de gestionar la planificación de las tareas mediante el uso de issues y milestones.

En el siguiente enlace se puede encontrar el repositorio del trabajo: \url{https://github.com/smi0010/TFG_Herramientas_Realidad_Aumentada}.

Esta herramienta ha sido escogida por tener experiencia trabajando con ella en asignatura previas del grado.

\subsubsection{ZenHub}\label{ZenHub}
ZenHub\footnote{\url{https://www.zenhub.com/}} se trata de una herramienta para la gestión de proyectos. Se encuentra disponible para GitHub mediante una extensión del navegador.

ZenHub permite usar un tablero Kanban para el seguimiento y organización del proyecto. En el se representan las issues como tarjetas, que las podemos organizar en diferentes columnas.


\subsubsection{GitHub Desktop}
 GitHub Desktop\footnote{\url{https://desktop.github.com/}}, se trata de un software  que facilita el uso del sistema Git mediante una interfaz gráfica, permitiendo de esta forma uso más sencillo y intuitivo. Permite sincronizarse con los repositorios de GitHub, y escoger cual usar en cada momento de una manera  rápida y sencilla.
 
 
 Actualmente se encuentra disponible para sistemas Windows y masOS.

\subsection{Herramientas de desarrollo}

\subsubsection{Unity}
Unity\footnote{\url{https://unity.com/es}}, es un motor gráfico de desarrollo de videojuegos creado por Unity technologies, disponible para Windows, Linux y Mac OS. En él, se pueden desarrollar aplicaciones para distintas plataformas como, Windows, MacOS, Linux, Andorid, iOS, videoconsolas, WebGL, tvOS y Facebook.

\subsubsection{VisualStudio}
ViusalStudio\footnote{\url{https://visualstudio.microsoft.com/es/}}, se trata de un entorno de desarrollo para Windows, Linux y macOS. Es compatible con múltiples lenguajes de programación. C\# es uno de los principales lenguajes con los que es compatible, y dado que los scripts de Unity son en dicho lenguaje, se ha escogido visual studio para su desarrollo.


\subsubsection{Vuforia}
Como se ha explicado en el apartado \ref{sub:Def_Vuforia}, Vuforia\footnote{\url{https://developer.vuforia.com/}} se trata de un SDK para aplicaciones de AR, para las plataformas de Android, iOS, Windows, Unity y HoloLens. Proporciona las herramientas necesarias para desarrollar realidad aumentada en el entorno Unity.

\subsubsection{CoSpace}
Como se expuso en el apartado \ref{sub:Def_Vuforia}, CoSpaces se trata de una plataforma de apoyo educativo. Se puede acceder a ella tanto por web, como desde la aplicación de móvil. Permite crear entornos virtuales y de realidad aumentada para el MergeCube, desde su entorno de desarrollo web.

\subsubsection{Gimp}
Gimp\footnote{\url{https://www.gimp.org/}}, se trata de un software libre y gratuito, de edición de imágenes digitales. Se ha usado para la creación y edición de algunos de los sprites usados en la aplicación.

\subsubsection{Pencil}
Pencil\footnote{\url{https://pencil.evolus.vn/}} se trata de un software gratuito para el diseño de prototipos de interfaces gráficas para aplicaciones. Esta pensada para el diseño de interfaces tanto para dispositivos móviles como aplicaciones de escritorio.

\section{Técnicas utilizadas en el desarrollo del trabajo}
\subsubsection{Scrum}
Scrum se trata de una metodología de desarrollo ágil de software que se encuentra basada en el manifiesto ágil~\footnote{\url{https://agilemanifesto.org/principles.html}}. Se ha seguido esta metodología realizando srpints con una duración normalmente de una semana, y una reunión semanal al finalizar cada sprint en las que se mostraban los avances realizados y se planteaba la planificación del siguiente.

Para facilitar el seguimiento de esta metodología se ha utilizado en Github(\ref{github}) las issues y la extensión de ZenHub(\ref{ZenHub}), que permite organizar de una forma mas cómoda mediante tableros kanban\footnote{\url{https://es.wikipedia.org/wiki/Kanban_(desarrollo)}}.


REVISADO HASTA AQUI DE MOMENTO; AHORA ESTOY SIGO TRABAJANDO EN EL RESTO DE APARTADOS


\capitulo{5}{Aspectos relevantes del desarrollo del proyecto}
En este apartado se explicarán los aspectos mas relevantes del desarrollo del proyecto. Las decisiones que se han tomado a lo largo del proyecto, los problemas que han surgido.

\subsection{Formación}
Para este proyecto era necesario aumentar conocimientos de los que se carecían. Principalmente en  Unity y las herramientas de realidad aumentada. Pues los conocimientos que tenía de antemano sobre Unity no eran muy grandes y aun así era necesario actualizarlos de cara al proyecto.

Para ello, se tenían como referencia los recursos oficiales del \href{https://docs.unity3d.com/2019.2/Documentation/Manual/UnityManual.html}{manual de Unity} y sus foros de ayuda de la comunidad. Del mismo modo, se han seguido los recursos oficiales del \href{https://library.vuforia.com/getting-started/overview.html}{manual de Vuforia}. Estos recursos también se han completado con diferentes videotutoriales, tanto creados por las propias compañías como por parte de la comunidad.

Una de las ventajas con Unity, era el uso de C\# pues era un lenguaje en el que ya tenía algunos conocimientos, lo cual ha facilitado el proceso de desarrollo en esta parte. 

\subsection{Estudio herramientas}
Uno de los principales puntos de este proyectos ha sido el estudio de diferentes herramientas de realidad aumentada. Para la selección de las herramientas que se estudiarían se siguieron varios pasos. 
Las primeras seleccionadas fueron ARCore y ARKit, pues que gracias al nombre de sus compañías ya había oído hablar de ellas incluso antes de empezar a investigar. Otra de las primeras herramientas se trataba de CoSpaces y el MergeCube, pues desde el comienzo del proyecto uno de los tutores le parecía interesante especialmente por su aplicación en docencia. Lamentablemente el uso exclusivo del SDK del MergeCube no fue posible debido a que se dejo de dar soporte de desarrollo, por lo que ya no era posible descargar el SDK ni conseguir las claves de licencia para su uso.

Para el resto de herramientas seleccionadas, se procedió a realizar una investigación sobre las herramientas que actualmente son mas populares y aquellas que destaquen por sus funcionalidades.

\subsection{Elección de Herramienta de realidad aumentada para el caso práctico}

Uno de los objetivos del proyecto era el de realizar una aplicación en la que poder poner a prueba las funcionalidades de la realidad aumenta. Para ello se debió de escoger entre una de las herramientas estudiadas. para esta selección se han tenido varios aspectos en cuenta: 
\begin{itemize}
	\item La capacidad de esa herramienta para funcionar en el mayor número de dispositivos posibles.
	\item Que se pueda desarrollar una aplicación de forma gratuita.
	\item La continuidad de dicha herramienta, es decir que esa herramienta siga actualizándose y teniendo un soporte oficial.
	\item La comunidad que tiene. Con este aspecto me refiero a la cantidad de usuarios que usan la herramienta. Pues cuanto más grande sea una comunidad de usuarios en un software, nos puede facilitar el resolver futuros problemas.
	
\end{itemize}

ArCore y Arkit, dada su tecnología eran una de las primeras opciones, pero en el caso de ArKit unicamente esta disponible para dispositivos iOS. Y en el caso de ArCore el número de dispositivos Android compatibles aún es algo limitado.
En el caso de las aplicaciones de MetaVerse, ZapWorks, WikiTude y 8thWall, se han descartado dado que para para el desarrollo de una aplicación únicamente ofrecen una versión gratuita por un tiempo limitado.

En el caso de OpenCv, su principal problema es que en comparación con otras herramientas de realidad aumentada, OpenCv se queda limitada en cuanto a sus posibilidades.

Finalmente la decisión estaba entre Kudan y Vuforia. Con kudan que parece ser un software muy prometedor, el problema estaba en que el SDK disponible no fui capaz de funcionara, es probable que se deba a que lleve tiempo sin actualizarse, pues la ultima versión de Unity en la que esta testado era la 2018. sumado ha este problema esta también el que no tiene una comunidad muy grande ni una guía muy detallada de uso, por lo que buscar posibles soluciones a dudas o problemas se complicaría bastante.
Es por esto que finalmente me decante por Vuforia, pues  ofrece la posibilidad de desarrollar gratuitamente, siempre que no se publique con el objetivo comercial. Es compatible 


\subsection{Elección de Unity}
Para desarrollar un ejemplo que aplicara las herramientas estudiadas, se tuvo en cuenta aquellos entornos con los que las herramientas eran compatibles, ademas de a ser posible buscar un entorno gratuito. 

La mayoría de las herramientas contaban con soporte oficial para el desarrollo en Unity. Y dado que la aplicación a desarrollar era de realidad aumentada Unity era la mas apropiada, pues para el manejo de modelos 3D y otros de estos recursos, era el entorno que mas facilidades otorgaba. Además de que cuenta con um amplio manual de usuario, y un gran apoyo de la comunidad.

\subsection{Problemas Instalación Vuforia en Unity}
Para incorporar en el proyecto de Unity la herramienta de Vuforia, hay varios métodos, en este caso yo decidí hacerlo desde la Asset Store de Unity. Para poder instalarlo, simplemente debemos buscar en la tienda Vuforia. Entre los resultados podemos encontrar diferentes Assest que ofrece PTC, la compañía de Vuforia. 

Hasta aquí todo bien, pero en vez de seleccionar Vuforía Engine, que contiene unicamente el engine de Vuforia, seleccione y importe Vuforia Core Samples, que se trata del engine de Vuforia pero viene con múltiples ejemplos desarrollados de Vuforia. De este hecho no me percaté hasta bastante más adelante. 

Me pude dar cuenta de esto por el peso de la aplicación y es que el APK era demasiado pesado para el contenido que estaba utilizando. A ese elevado peso de la APK, no le di demasiada importancia hasta que hice pruebas en un proyecto aparte, en los que solo instale Vuforia Engine y pude comprobar que en estos proyectos el peso era considerablemente menor.

Ademas del peso de la aplicación, no suponía ningún problema en el funcionamiento, pues realmente el engine era el mismo solo que en uno tenía mas recursos. Dado que a la larga el este excesivo peso aun que no fuera un problema, si sería un inconveniente. Empecé a revisar el proyecto para borrar los elementos que hacían que el peso de la aplicación ascendiera tanto. Finalmente si les puede encontrar en la carpeta \textbf{StreamingAssets}, \textbf{Vuforia}, y borrar todos los elementos excepto el que tendrá el nombre del prefab del MergeCube añadido.

\subsection{Corrupción del proyecto}

En dos ocasiones el proyecto de Unity se corrompió, de forma que la única opción fue restaurarlo con una copia de seguridad. 

En ambas ocasiones el problema fue el mismo, todos los scripts pasaban a ser inutilizables en el proyecto. Esto impedia la posibilidad de asociar un script a un GameObject para que este usara las funciones del script. Este problema no solo impedía realizar nuevas asociaciones, sino todos los objetos que ya estaban asociados a algún script ya no lo estaban.

La lamentablemente no puede encontrar el origen del problema, ni una solución. Algunas de las posibles soluciones que encontré en foros, eran borrar los ficheros temporales o reimportar el proyecto. La lamentablemente ninguno funciono y tuve que restaurarlo con un copia. 





\capitulo{6}{Trabajos relacionados}

Son numerosos los trabajos y estudios que hoy en día se realizan respecto a las tecnologías de realidad aumentada, en claro reflejo de su creciente interes....  En este apartado veremos algunos de los trabajos que destacan.

\section{Artículos}
\subsection{Realidad Aumentada, Educación y Museos}
 
 Se trata de un estudio sobre la realidad aumentada

\section{Libros}

\section{Paginas Web}


\subsection{Niantic}
Niantic~\footnote{\url{https://nianticlabs.com/es/}} se trata de una empresa centrada en el desarrollo de aplicaciones de realidad aumentada. Actualmente sus aplicaciones más conocidas son Ingress, Pokemon Go y HarryPotter wizzards unite.
Pero aparte de estas aplicaciones también se dedican a la investigación y desarrollo de la realidad aumentada.

Uno de los destacados es del que hablan en la entrada <<Building experiences from the ground up>> de su blog, en el que explican su intención de crear mejores experiencias de realidad aumentada con
la capacidad de reconocer objetos del entorno para crear mejores sensaciones de profundidad. Para ello es necesario transformar una imagen 2D en un espacio 3D. Un ejemplo sería, en una imagen en la que salen varios arboles, si queremos poner un elemento de realidad aumentada, al moverse ese elemento no distinguiría el fondo de un árbol, por lo que la imagen AR siempre estará en un plano por encima. Pero si la imagen 2D pasa a ser un mapa 3D en el que se diferencian los arboles del fondo, sería posible que ese elemento de realidad aumentada pase de una forma realista detrás de los arboles. 
\imagen{Trabajos_Relacionados/nianticAR}{ Representación del reconocimiento de profundidad.~\cite{nianticAR}}
\url{https://nianticlabs.com/es/blog/building-experiences-from-the-ground-up/}



\subsection{-}
\capitulo{7}{Conclusiones y Líneas de trabajo futuras}


\section{Conclusiones}
A pocos días de la entrega, creo poder decir que si se han cumplido los objetivos plantados al comienzo del proyecto. Se ha realizado un estudio de las diferentes herramientas de realidad aumentada actuales, y se ha conseguido realizar una aplicación que pone aprueba las funcionalidades de una de estas herramientas, pensando en un ámbito educativo. Lamentablemente no ha sido posible implementar todas las funcionalidades de realidad aumentada que nos habría gustado, pues algunas de estas funcionalidades únicamente están disponible en ciertos dispositivos móviles.

Se han adquirido conocimientos del desarrollo en Unity y de algunas de sus extensiones. También sobre las herramientas de realidad aumentada y que aspectos tener en cuenta a la hora de escoger una herramienta.

El desarrollo de la aplicación al principio ha sido lento, pues en las primeras semanas no teníamos una idea clara sobre que tipo de ejercicio educativo realizar. Y es que algunas ideas, aun que eran buenas, su complejidad a la larga habría resultado problemática. Finalmente nos decantamos por simular las fases del crecimiento de una planta, pues dados los conocimientos que tenía me pareció que era una idea asumible de hacer.

Uno de los problemas que surgieron a lo largo de la aplicación, fue a la hora escoger modelos 3D para utilizar. Si bien es cierto que hay multitud de modelos gratuitos, no siempre tienes la surte de que el modelo que mejor se puede acoplar a tu proyecto, pertenezca a uno de esos pack gratuitos. Por ejemplo a la hora de buscar plantas encontrar modelos para representar varias de sus fases es algo complejo, para solventar esto se decidió que varias de las fases fuesen con el mismo modelo, pero cambiando su escala.

El último problema ha destacar de la aplicación es el de la falta de experiencia con Unity, aunque si había trabajado con él anteriormente, no había realizado nada particularmente complejo. Esto lo he notado a la hora de escoger, diseñar o configurar algunos elementos. Por ejemplo, tuve problemas al inicio para conseguir detectar las partículas de agua que colisionaban con la planta.

Considero que uno de las aspectos más importantes que he aprendido con este proyecto ha sido a gestionar mi tiempo de una manera mas eficiente.

Aunque que considero que la aplicación podría ser mejorable y que tuviera aun algo mas de contenido, como podría haber sido otro ejemplo aparte del de las plantas. Si estoy satisfecho con el resultado, y especialmente con los conocimientos adquiridos de las herramientas usadas y de las investigaciones realizadas.

\section{Lineas de trabajo futuras}
Tras concluir el desarrollo, puedo pensar en algunas tareas para posibles trabajos futuros.
\begin{itemize}
	\item En primer lugar, teniendo en cuenta uno de los aspectos que mas relevantes de la realidad aumentada hoy en día es la tendencia ha dejar de lado los marcadores físicos y centrarse en el reconocimiento sin marcadores, el reconocimiento del entorno que nos rodea.
	Por ello una de los aspectos en los que se podría mejorar la aplicación es implementado este tipo de realidad aumentada, de una forma mas precisa y con mayor papel en el uso de la app.
	
	Si bien es cierto, que la principal limitación de esta realidad aumenta se encuentra en los dispositivos compatibles con ella, pero bastante posible que en unos años esta limitación pueda ser superada sin problemas.
	\item Llevar la aplicación a otras plataformas como iOS o Windows. Por ejemplo sería interesante llevar la aplicación a Windows, usando Webcam.
	\item Uso de la realidad aumentada Web. Especialmente con las mejoras en las lineas de telecomunicaciones con la conexión 5G, permitirá realizar grandes avances en la realidad aumentada.
	\item Uso de otras herramientas de realidad aumentada, para completar con nuevas funcionalidades o simplemente mejorar las ya existentes.
	\item Añadir la posibilidad de usar la realidad aumentada, junto a unas gafas VR (las gafas que se usan con el móvil), para poder experimentar la realidad aumentada con las manos libres. 
	
\end{itemize}




\bibliographystyle{plain}
\bibliography{bibliografia}

\end{document}
